% ---------------------------------------------------------------------
% Section 6
% ---------------------------------------------------------------------

\begin{cnread}
本节处理「加速电荷为什么会辐射」。\AEG{} 的答案出奇地干净：匀速运动停留在
固定核坐标 $b$ 的稳定子轨道上，绕任何闭回路的乘性和乐**望远镜抵消为 $1$**
（这就是「无辐射定理」）；一旦发生加速，核坐标（涟漪极点）跳动
$b_1\to b_2\ne b_1$，四步闭环上的交叉比和乐
$\Lambda(\square)=\dfrac{(1+a_2b_1)(1+a_1b_2)}{(1+a_1b_1)(1+a_2b_2)}$ 就不再等于
$1$——而且它既不依赖标架截面，也无法被望远镜抵消。本节还给出 $\mathbb F_7$
上的精确见证 $\Lambda=4$，并说明这正是 Paper~IV 定义 4.10 / 定理 4.11 所
记录的对象。
\end{cnread}

\subsection{Ripple pencils with a moving tangency point}

Paper~0, Subsection~4.3, records the geometric state that becomes available once
the second slot is read projectively:

\begin{quote}\small
``Thus a general right-expansion history should not be pictured as one fixed
family of concentric circles.  \emph{Each prefix has its own ripple pencil, and
the common tangency point of that pencil may move with $k$.}  The appropriate
geometric state is $(G_k,G_k(z_0),p_k,\mathcal R_k)$: the operator, the selected
point, the pole, and the pulled-back wavefront family.''
\end{quote}

The four-tuple is the whole content of the radiation phenomenon in this
framework: the operator carries the cumulative transformation, the selected
point carries the current value, the pole carries where the wavefronts are
tangent, and the pencil carries the wavefronts themselves.  A history in which
$p_k$ is constant emits nothing new; a history in which $p_k$ jumps emits a
mismatch between two pencils.

\subsection{Bivaluations and the comoving frame}

Fix $P_k=\begin{pmatrix}A_k&B_k\\C_k&D_k\end{pmatrix}\in\PGL_2(\R)$ with
$\Delta_k=A_kD_k-B_kC_k>0$, acting on the reference pair
$(q_0^+,q_0^-)=(0,\infty)$.  The image and kernel directions are
\begin{equation}\label{eq:bivaluation}
  q_k^+=P_k(0)=\frac{B_k}{D_k}=-a_k,
  \qquad
  q_k^-=P_k(\infty)=\frac{A_k}{C_k}=\frac1{b_k},
\end{equation}
and the pole of the contravariant pullback is
$p_k=P_k^{-1}(\infty)=-D_k/C_k$ (Paper~0, Theorem~4.2 and Section~5.1).  The
pair $(a_k,b_k)$ is a regular bivaluation in the sense of Paper~IV,
Definition~3.1, provided $1+a_kb_k\ne0$.

\begin{theorem}[Bivaluation--projector correspondence, {\cite[Thm.~3.2]{AEG4}}]
\label{thm:projector}
There is a natural bijection between regular bivaluations and rank-one
idempotents of $K^{2}$.  In the affine chart with
$v_a=\binom{-a}{1}$ and $\varphi_b=(-b\;\;1)$,
\begin{equation}\label{eq:projector}
  \Pi_{a,b}=\frac{v_a\otimes\varphi_b}{\varphi_b(v_a)}
  =\frac{1}{1+ab}\begin{pmatrix}ab&-a\\-b&1\end{pmatrix},
  \qquad
  \varphi_b(v_a)=1+ab .
\end{equation}
\end{theorem}

The projector $\Pi_k=\Pi_{a_k,b_k}$ is the intrinsic object behind the
``comoving frame'' of the physical reading: it splits the value space into the
image line $L_k^+$ (the longitudinal near field) and the kernel line $L_k^-$ (the
transverse wavefront direction), and it is the unique such splitting determined
by the history's prefix.  Everything in this section is a statement about
products of these projectors.

\subsection{Uniform motion: no radiation}

\begin{theorem}[Fixed-kernel flatness]
\label{thm:no-radiation}
Let a history have constant kernel coordinate $b_k\equiv b_0$ for all prefixes,
and let the image coordinate return to its initial value,
$a_m=a_0$.  Then the multiplicative holonomy
\begin{equation}\label{eq:lambda-edge}
  \lambda(e_k):=\frac{1+a_kb_{k-1}}{1+a_{k-1}b_{k-1}}
\end{equation}
satisfies $\prod_{k=1}^{m}\lambda(e_k)=1$.  In particular no closed history with
constant kernel coordinate carries a non-unit multiplicative holonomy.
\end{theorem}

\begin{proof}
With $b_{k-1}=b_0$ the product is
\[
  \prod_{k=1}^{m}\frac{1+a_kb_0}{1+a_{k-1}b_0}
  =\frac{1+a_1b_0}{1+a_0b_0}\cdot\frac{1+a_2b_0}{1+a_1b_0}\cdots
   \frac{1+a_0b_0}{1+a_{m-1}b_0}=1,
\]
each numerator cancelling the next denominator, the last against the first
because $a_m=a_0$.
\end{proof}

The theorem is the algebraic content of the second reading of
Section~5.4: uniform motion keeps the ripple pencil's tangency point fixed, and
a fixed tangency point gives no holonomy.  By
Hypothesis~\ref{hyp:radiation} this is the no-radiation statement.

\subsection{Accelerated motion: the pole jumps}

Now let the history include the inverting element $J(z)=-1/z$, i.e. a second-slot
division.  Then the prefix leaves the affine Borel sector and enters the other
Bruhat cell, and the kernel coordinate becomes nonzero.  A physical acceleration
is a change of the kernel coordinate between one step and the next.

\begin{theorem}[Cross-ratio holonomy]
\label{thm:cross-ratio}
Consider the closed four-step base loop
\[
  \square:\quad (a_1,b_1)\to(a_2,b_1)\to(a_2,b_2)\to(a_1,b_2)\to(a_1,b_1),
\]
in which the image coordinate moves with the kernel held fixed, then the kernel
moves with the image held fixed, and then both return.  Assume every vertex and
every step is transverse, i.e. all of $1+a_ib_j$ and the two cross terms
$1+a_2b_1$, $1+a_1b_2$ are nonzero.  Then the multiplicative holonomy of the
loop is
\begin{equation}\label{eq:cross-ratio}
  \Lambda(\square)
  =\frac{(1+a_2b_1)(1+a_1b_2)}{(1+a_1b_1)(1+a_2b_2)}
  =1-\frac{(a_2-a_1)(b_2-b_1)}{(1+a_1b_1)(1+a_2b_2)} .
\end{equation}
It is invariant under changes of the frame lifts $\Phi_k$, and
\begin{equation}\label{eq:cross-ratio-nontrivial}
  \Lambda(\square)\ne1
  \quad\Longleftrightarrow\quad
  a_2\ne a_1\ \text{and}\ b_2\ne b_1 .
\end{equation}
Moreover, as an identity between projectors,
\begin{equation}\label{eq:projector-product}
  \Pi_{a_1,b_1}\Pi_{a_1,b_2}\Pi_{a_2,b_2}\Pi_{a_2,b_1}
  =\Lambda(\square)^{-1}\,\Pi_{a_1,b_1},
\end{equation}
so $\Lambda(\square)$ is a gauge-invariant observable of the four idempotents
alone.
\end{theorem}

\begin{proof}
Along $e_1:(a_1,b_1)\to(a_2,b_1)$ and $e_3:(a_2,b_2)\to(a_1,b_2)$ the factor
\eqref{eq:lambda-edge} equals $(1+a_2b_1)/(1+a_1b_1)$ and
$(1+a_1b_2)/(1+a_2b_2)$ respectively, because the kernel coordinate is the
\emph{previous} one in each case.  Along $e_2$ and $e_4$ the image coordinate is
constant, so $(1+a_kb_{k-1})/(1+a_{k-1}b_{k-1})=1$.  The product is
\eqref{eq:cross-ratio}.  Gauge invariance: a change $\Phi_k\mapsto c_k\Phi_k$
multiplies each $\lambda(e_k)$ by $c_k/c_{k-1}$, whose product over a closed loop
is $1$.  For \eqref{eq:cross-ratio-nontrivial}, subtract numerator from
denominator of \eqref{eq:cross-ratio}: the difference is
$(a_2-a_1)(b_2-b_1)$, which is nonzero exactly when both coordinates move.
Identity \eqref{eq:projector-product} is a direct product of the rank-one
matrices \eqref{eq:projector}; it is proved symbolically in
Appendix~\ref{app:matrices} and checked exhaustively over $\F_7$ there.
\end{proof}

\begin{corollary}[Larmor criterion, algebraised]
\label{cor:larmor}
A closed history radiates, in the sense of Hypothesis~\ref{hyp:radiation},
precisely when both the image coordinate and the kernel coordinate move.  In the
infinitesimal limit the log-holonomy is
\begin{equation}\label{eq:larmor}
  \log\Lambda(\square)
  =-\frac{(a_2-a_1)(b_2-b_1)}{(1+a_1b_1)(1+a_2b_2)}+O(3)
  =-\iint_{\square}\frac{\mathrm da\wedge\mathrm db}{(1+ab)^{2}}+O(3),
\end{equation}
so that displacement alone and velocity change alone are both non-radiative,
while their product is radiative.
\end{corollary}

\begin{proof}
The first equality is the logarithm of \eqref{eq:cross-ratio} expanded with
$(a_2-a_1)(b_2-b_1)$ small.  The second is the same expansion written as an
integral over the coordinate rectangle; Check~V3.1 verifies it symbolically.
The vanishing statement is \eqref{eq:cross-ratio-nontrivial}.
\end{proof}

\begin{remark}[The structure of the answer]
\label{rem:larmor-structure}
The criterion has the shape of the classical radiation statement --- a charge at
rest radiates nothing, a uniformly moving charge radiates nothing, and an
accelerated one does --- but the two ``nothing'' cases are proved by different
mechanisms: the rest case has $da=0$ and the uniform case has $db=0$.  The
criterion is therefore not ``velocity change causes radiation'' but
``\emph{simultaneous} displacement and change of the wavefront tangency point
causes radiation'', which in the affine sector reduces to the classical
statement.
\end{remark}

\begin{figure}[t]
\centering
\begin{tikzpicture}[x=1.05cm,y=1.05cm,>=Latex,every node/.style={font=\scriptsize}]
  \fill[blue!3] (-2.7,-0.7) rectangle (7.5,3.9);
  \draw[->,black!45] (-1.0,0)--(5.4,0) node[right] {$a$};
  \draw[->,black!45] (0,-0.4)--(0,3.5) node[above] {$b$};
  \coordinate (v1) at (0.6,0.55);
  \coordinate (v2) at (3.6,0.55);
  \coordinate (v3) at (3.6,2.75);
  \coordinate (v4) at (0.6,2.75);
  \draw[very thick,black!70] (v1)--(v2);
  \draw[very thick,blue!60!black] (v2)--(v3);
  \draw[very thick,black!70] (v3)--(v4);
  \draw[very thick,blue!60!black] (v4)--(v1);
  \foreach \p/\lab/\pos in {v1/{(a_1,b_1)}/below left,
                            v2/{(a_2,b_1)}/below right,
                            v3/{(a_2,b_2)}/above right,
                            v4/{(a_1,b_2)}/above left}
    {\fill (\p) circle (3pt) node[\pos] {$\lab$};}
  \node[black!70,below] at (2.1,0.55) {image moves, $b$ fixed};
  \node[black!70,above] at (2.1,2.75) {image returns, $b$ fixed};
  \node[blue!60!black,align=left,anchor=west,text width=2.5cm] at (4.05,1.65)
    {kernel moves at fixed $a$: this step alone has $\lambda=1$};
  \node[blue!60!black,align=right,anchor=east] at (-0.4,1.65)
    {kernel returns:\\ $\lambda=1$};
  \node[align=center] at (2.2,-1.35)
    {$\Lambda(\square)=\dfrac{(1+a_2b_1)(1+a_1b_2)}{(1+a_1b_1)(1+a_2b_2)}\neq1$
     \ iff both $a$ and $b$ move};
\end{tikzpicture}
\caption[The four-step bivaluation square]%
{The four-step bivaluation square.  Each individual step has a corrected base
factor, but only the two image-moving steps contribute; the product is the
cross-ratio \eqref{eq:cross-ratio}.  A loop in which the kernel coordinate never
moves collapses to $\Lambda=1$ by Theorem~\ref{thm:no-radiation}.}
\label{fig:square}
\end{figure}

\subsection{An exact finite witness}

\begin{example}[$\F_7$ calibration]
\label{ex:f7}
Let $K=\F_7$, so that $|G|=336$, $|\Biv|=56$ and $|H|=6$ by Paper~IV,
Proposition~3.7 and Example~3.8.  Take the loop
\[
  (a_1,b_1)=(0,0),\quad (a_2,b_1)=(1,0),\quad (a_2,b_2)=(1,1),\quad (a_1,b_2)=(0,1).
\]
All four vertices are transverse, since $1+a_ib_j\in\{1,2\}\subset\F_7^\times$.
Then
\[
  \Lambda(\square)\equiv\frac{(1+0)(1+0)}{(1+0)(1+1)}\equiv 2^{-1}\equiv 4 \pmod 7,
\]
so $\Lambda(\square)=4\ne1$ in $H\cong\F_7^{\times}$.
\end{example}

This is the exact finite witness that the obstruction is not an artefact of the
real or complex topology: the holonomy is computed in a finite group of order
$336$, and it is a nonidentity element of the fiber $H$.  Appendix~\ref{app:matrices}
records the same example together with the exhaustive check of
\eqref{eq:projector-product} over all transverse loops in $\F_7$.

\subsection{The physical dictionary, stated as a dictionary}

\begin{warning}[What Hypothesis~\ref{hyp:radiation} adds]
\label{warn:radiation}
Theorem~\ref{thm:no-radiation} and Theorem~\ref{thm:cross-ratio} are theorems
about idempotents.  They say: a closed history whose kernel coordinate is
constant has a flat multiplicative transport, and a closed history that moves
both coordinates has a holonomy with an exact cross-ratio value.  They do not by
themselves say that anything is emitted.  The identification of
$\Lambda(\square)\ne1$ with radiation is Hypothesis~\ref{hyp:radiation}, and the
identification of its magnitude with radiated power is a further, unproved
refinement.

Three dictionary readings are available, and each should be judged on its own:
\begin{enumerate}[label=(\alph*)]
  \item \emph{Wavefront dislocation.}  The two pencils $\mathcal R_k$ and
  $\mathcal R_{k+1}$ are tangent at different ideal points, so at their interface
  the wavefronts cannot nest smoothly; the mismatch is a caustic, i.e. a kink in
  the propagating wavefront.  The magnitude of the mismatch is
  $|p_{k+1}-p_k|$ together with the holonomy $\Lambda$.
  \item \emph{Field falloff.}  In electrodynamics the field of an accelerated
  source has a term falling off as $1/r$ and proportional to the transverse
  acceleration, while the fields of static and uniformly moving sources fall off
  faster.  The correspondence asserted by Hypothesis~\ref{hyp:radiation} is that
  the non-telescoping holonomy is the \AEG{} carrier of that $1/r$ term.  No
  fall-off exponent is derived here.
  \item \emph{Phase residue.}  By Remark~\ref{rem:prequantum} the loop's
  non-closure along the Reeb fiber measures a phase; a phase that does not return
  to itself is the same kind of statement as a radiated quantum's escape.
\end{enumerate}
\end{warning}

\subsection{Relation to the \AEG{} manuscripts}

The construction proved in this section is not new to the \AEG{} series: it is
Paper~IV, Definition~4.10 and Theorem~4.11, with the $\F_7$ witness of
Example~4.13.  It was developed in the dialogue summarised here and recorded in
the manuscript under the heading of the projective-torsion problem.  The present
section's contribution is therefore not the theorem but the reading: the
non-telescoping holonomy of Paper~IV, Theorem~4.11 is the \AEG{} object
that the physical narrative calls ``radiation'', and
Corollary~\ref{cor:larmor} states the Larmor criterion in its terms.

Section~\ref{sec:torsion} completes the picture by packaging the same data into a
transport that also reproduces the affine torsion of Paper~0 on the Borel
sector.
